\documentclass[11pt, a4paper]{report}
\usepackage[margin=0.8in]{geometry}
\usepackage{amsmath, amssymb, amsthm}
\usepackage{tcolorbox}
\usepackage[colorlinks=true, linkcolor=blue, urlcolor=blue, citecolor=blue, bookmarks=true, bookmarksopen=true]{hyperref}

\setlength{\parindent}{0pt}
\setlength{\parskip}{0.5em}

\newtheorem{theorem}{Theorem}

\title{\textbf{Group Theory Proof Verification}\\ \Large Graduation Project Testing}
\author{General Relativity Graduation Project}
\date{\today}

\begin{document}

\maketitle
\tableofcontents
\newpage

\chapter{Alkali}
Formulas identified:Wavenumber of the Sharp SeriesFormula: $\bar{\nu}_{sharp} = R \left[ \frac{1}{(n_0 - \Delta_p)^2} - \frac{1}{(n - \Delta_s)^2} \right]$Meaning of each variable: $\bar{\nu}_{sharp}$ is the wavenumber of the emitted photon. $R$ is the Rydberg constant. $\Delta_p$ and $\Delta_s$ are the quantum defects for the $p$ ($l=1$) and $s$ ($l=0$) orbitals, respectively. $n_0$ is the lowest available principal quantum number for the $p$ state, and $n$ is the varying principal quantum number of the upper $s$ state.When it is used: To calculate the exact position of spectral lines when an electron falls from an excited $ns$ state to the lowest $np$ state.Common exam situations: Using the provided Quantum Defect Table (Table 5.1) to find $\bar{\nu}$ for a specific transition in sodium (e.g., $4s \to 3p$).Wavenumber of the Principal SeriesFormula: $\bar{\nu}_{principal} = R \left[ \frac{1}{(n_0 - \Delta_s)^2} - \frac{1}{(n - \Delta_p)^2} \right]$Meaning of each variable: Similar to the above, but the transition is from an excited $np$ state down to the fundamental $n_0 s$ ground state.When it is used: Used for the most prominent ("strong") absorption and emission lines in alkali metals.Common exam situations: Calculating the convergence limit of the principal series (by letting $n \to \infty$).Wavenumber of the Diffuse SeriesFormula: $\bar{\nu}_{diffuse} = R \left[ \frac{1}{(n_0 - \Delta_p)^2} - \frac{1}{(n - \Delta_d)^2} \right]$Meaning of each variable: Involves transitions from $nd$ states down to $n_0 p$ states, utilizing $\Delta_p$ ($l=1$) and $\Delta_d$ ($l=2$).When it is used: To find the wavenumbers of the narrow, compound triplet lines in the diffuse series.Common exam situations: Comparing energy levels of different angular momentum states using quantum defect values.Wavenumber of the Fundamental SeriesFormula: $\bar{\nu}_{fundamental} = R \left[ \frac{1}{(n_0 - \Delta_d)^2} - \frac{1}{(n - \Delta_f)^2} \right]$Meaning of each variable: Involves transitions from $nf$ states down to $n_0 d$ states.When it is used: To calculate lines lying primarily in the Infrared (IR) region.Common exam situations: Explaining why the quantum defect for $f$-orbitals is negligible (as seen in Table 5.1) making these transitions very similar to pure Hydrogen transitions.Rydberg-Schuster LawFormula: $\bar{\nu}_{P\infty} - \bar{\nu}_{S\infty} = \bar{\nu}_{P1}$Meaning of each variable: $\bar{\nu}_{P\infty}$ and $\bar{\nu}_{S\infty}$ are the convergence limits (wavenumbers at $n \to \infty$) of the Principal and Sharp series, respectively. $\bar{\nu}_{P1}$ is the wavenumber of the first line of the principal series.When it is used: To relate the ionization limits of different orbital series to one another.Common exam situations: Given the limit of the Sharp series and the first Principal line, calculate the Principal series limit.Runge's LawFormula: $\bar{\nu}_{S\infty} - \bar{\nu}_{F\infty} = \bar{\nu}_{d1}$ (Text notes the Sharp/Diffuse limits are identical: $\bar{\nu}_{S\infty} = \bar{\nu}_{d\infty}$)Meaning of each variable: Relates the Sharp/Diffuse series limits to the Fundamental limit via the first line of the diffuse series.When it is used: Establishing energy differences between the lowest $p$, $d$, and $f$ states.Common exam situations: Proving the internal consistency of the term values.Important exam targets:Mastering Selection Rules: You must be able to strictly apply $\Delta l = \pm 1$, $\Delta j = 0, \pm 1$, and $\Delta m_j = 0, \pm 1$. The exam will likely ask you to "derive" the number of spectral lines for a given magnetic field (Zeeman) or spin-orbit coupling (Fine Structure) exactly as done for the D1 and D2 lines.Reading Grotrian Diagrams: Understand how to read Figure 5.1 (Lithium) and Figure 5.2 (Sodium) to identify the initial and final states of the Principal, Sharp, Diffuse, and Fundamental series.Quantum Defects: Be prepared to explain why the quantum defect is large for 's' electrons (penetration of the core) and decreases for 'p', 'd', and 'f' electrons based on the data in Table 5.1.Compound Doublets/Triplets: You will likely need to explain why the Sharp and Principal series are simple doublets, while the Diffuse and Fundamental series are compound triplets (due to the fine structure splitting of both the upper and lower states).

\chapter{LEC 26}
Important exam targets:The Born-Oppenheimer Justification: You must be able to state why we can separate nuclear and electronic wavefunctions (the massive difference in mass and speed).Energy Balance: Be prepared to explain the exact mechanism of the potential energy minimum (Coulomb attraction vs. Nuclear-nuclear repulsion).Orbital Splitting Rule: Understand that when degenerate atomic energy levels interact to form a bond, they split. The total number of resulting molecular orbitals is always exactly equal to the initial number of interacting atomic orbitals.$H_2^+$ Stability: Be able to explain that the $H_2^+$ ion is mathematically stable specifically because its bonding energy at the equilibrium inter-nuclear distance is strictly lower than the energy of an individual, isolated hydrogen atom.
\chapter{Lecture--Molecular Orbital Theory(edited).pdf}
Formulas identified:Separation of Atomic Orbital WavefunctionFormula: $\varphi_{xyz} = \varphi_{radial}(r) \cdot \varphi_{angular}(\theta, \phi) = R_{nl}(r) Y_{lm}(\theta, \phi)$Meaning of each variable: The 3D atomic orbital $\varphi_{xyz}$ is split into a radial component $R_{nl}$ (which depends on the principal $n$ and orbital $l$ quantum numbers and the distance from the nucleus $r$) and an angular component $Y_{lm}$ (Spherical Harmonics, which depend on $l$ and magnetic quantum number $m$, and the angles $\theta, \phi$).When it is used: This is the foundational mathematical definition of an atomic orbital before it undergoes linear combination (LCAO) to form molecular orbitals.


Formulas identified:Energy of the Non-Rigid RotorFormula: $\epsilon_J = E_J / hc = BJ(J+1) - DJ^2(J+1)^2 \text{ cm}^{-1}$Meaning of each variable: $\epsilon_J$ is the rotational energy in wavenumbers. $B$ is the rotational constant. $D$ is the centrifugal distortion constant. $J$ is the rotational quantum number.When it is used: To find the precise energy levels of a realistic rotating molecule where the bond stretches at high rotational speeds.Common exam situations: Comparing this formula to the rigid rotor to explain why real molecular energy levels are slightly lower than predicted by the rigid model.Centrifugal Distortion Constant ($D$)Formula 1 (Fundamental): $D = \frac{h^3}{32\pi^4 I^2 r^2 kc} \text{ cm}^{-1}$Formula 2 (Practical): $D = \frac{4B^3}{\bar{\nu}^2} \text{ cm}^{-1}$Meaning of each variable: $I$ is the moment of inertia, $r$ is bond length, $k$ is the force constant, $B$ is the rotational constant, and $\bar{\nu}$ is the fundamental vibrational frequency.When it is used: To relate a molecule's resistance to stretching (vibrational frequency) to its rotational distortion.Common exam situations: Using the practical formula ($D = 4B^3 / \bar{\nu}^2$) to calculate the fundamental vibrational frequency $\bar{\nu}$ if only rotational microwave data ($B$ and $D$) are known.Vib-Rotational Selection Rules (Branches)Formulas: - P branch: $\Delta v = +1, \Delta J = -1$Q branch: $\Delta v = +1, \Delta J = 0$ (Labeled "forbidden" in the notes for standard diatomics)R branch: $\Delta v = +1, \Delta J = +1$Meaning of each variable: $\Delta v$ is the change in vibrational state. $\Delta J$ is the change in rotational state.When it is used: To categorize the spectral lines that appear when a molecule simultaneously absorbs an infrared photon to jump to a higher vibrational state while also changing its rotational speed.Common exam situations: Identifying which lines correspond to which branch on a given spectrum diagram.Important exam targets:Analyzing the Vib-Rotational Energy Level Diagram: You must master the diagram on Page 3 (Image 7). Notice how the P-branch lines (green arrows) originate from higher $J''$ levels dropping to lower $J'$ levels ($\Delta J = -1$), causing them to appear at frequencies lower than the band center ($\nu_0$). Conversely, the R-branch lines (blue arrows) jump to higher $J'$ levels ($\Delta J = +1$), appearing at frequencies higher than $\nu_0$.Centrifugal Distortion Concept: Be prepared to explain physically why the spacing between rotational lines shrinks. (Answer from the notes: "The more quickly a diatomic molecule rotates the greater is the centrifugal force tending to move the atoms apart." This increases the moment of inertia $I$, which decreases $B$, thereby lowering the energy gaps).

\chapter{lecture28.pdf}
Formulas identified:Total Molecular HamiltonianFormula: $\left( -\frac{\hbar^2}{2\mu}\nabla_{nuclear}^2 - \frac{\hbar^2}{2m_e}\sum \nabla_i^2 + V \right)\Psi = E\Psi$Meaning of each variable: $\mu$ is the reduced mass of the two nuclei. $m_e$ is electron mass. $\nabla_{nuclear}^2$ and $\nabla_i^2$ are the Laplacian operators for the nuclei and the $i$-th electron, respectively. $V$ is the total Coulomb potential (nucleus-nucleus repulsion + electron-electron repulsion - electron-nucleus attraction).When it is used: As the starting point for all diatomic quantum mechanics derivations before applying approximations.Reduced MassFormula: $\mu = \frac{m_A m_B}{m_A + m_B}$Meaning of each variable: $m_A$ and $m_B$ are the masses of the two individual atomic nuclei.When it is used: To convert the two-body kinetic energy problem of the nuclei into a single-body problem orbiting the center of mass.Common exam situations: Calculating the rotational constant $B_e$ for specific isotopes (like HCl vs. DCl).Rotational Energy Levels (in Joules)Formula: $E_J = \frac{\hbar^2}{2I_e} J(J+1)$Meaning of each variable: $J$ is the rotational quantum number ($J = 0, 1, 2, ...$). $I_e = \mu R_e^2$ is the moment of inertia at equilibrium bond length.When it is used: To find the absolute quantized energy of a rotating molecule.Rotational Term Values (in Wavenumbers, $cm^{-1}$)Formula: $F(J) = \frac{E_J}{hc} = B_e J(J+1)$Formula for $B_e$: $B_e = \frac{h}{8\pi^2 c \mu R_e^2}$Meaning of each variable: $F(J)$ is the energy expressed in wavenumbers. $B_e$ is the rotational constant. $c$ is the speed of light.When it is used: Almost exclusively used for spectroscopic calculations, as spectrometers measure $cm^{-1}$ rather than Joules.Common exam situations: An exam problem will give you $B_e$ (or the spacing between spectral lines) and ask you to solve for the bond length $R_e$.Important exam targets:The Core Justification: You must be prepared to state exactly why the Born-Oppenheimer approximation works physically (the massive difference in inertia between nuclei and electrons) and how it manifests mathematically (dropping the derivative of the electronic wavefunction with respect to nuclear coordinates).Spectroscopic Units: Notice how $E_J$ is converted to $F(J)$. Exams will frequently try to trick you with units. Remember that $E = hc\bar{\nu}$, so dividing by $hc$ converts Joules to $cm^{-1}$.


Formulas identified:Rotational Term Value ($F(J)$)Formula: $F(J) = \frac{E_J}{hc} = B_e J(J+1)$Meaning of each variable: $F(J)$ is the rotational energy measured in wavenumbers ($cm^{-1}$). $B_e$ is the rotational constant. $J$ is the rotational quantum number ($0, 1, 2, ...$).When it is used: To map out the absolute energy positions of the rotating molecule.Common exam situations: Determining the population of a specific $J$ state relative to the ground state.Rotational Constant ($B_e$)Formula: $B_e = \frac{h}{8\pi^2 I_e c} = \frac{h}{8\pi^2 \mu R_e^2 c}$Meaning of each variable: $I_e$ is the moment of inertia. $\mu$ is the reduced mass. $R_e$ is the equilibrium bond length. $c$ is the speed of light.When it is used: To relate a macroscopic observable (spectral lines) to fundamental molecular geometry (bond length).Common exam situations: You are given the measured value of $B_e$ in $cm^{-1}$ and asked to isolate and calculate the exact bond length $R_e$ of a molecule (e.g., the HCl example on Page 8).Wavenumber of Rotational Spectral LinesFormula: $\bar{\nu} = F(J+1) - F(J) = 2B_e(J+1)$Meaning of each variable: $\bar{\nu}$ represents the exact position on a spectrum where an absorption line appears.When it is used: To predict the spacing between spectral lines.Common exam situations: Proving that the rigid rotor spectrum consists of equally spaced lines, with a constant gap of exactly $2B_e$ between each adjacent line.


Formulas identified:Harmonic Vibrational Term ValueFormula: $G(\upsilon) = (\upsilon + \frac{1}{2})\bar{\omega}_e$Meaning of each variable: $G(\upsilon)$ is the vibrational energy in wavenumbers. $\upsilon$ is the vibrational quantum number ($0, 1, 2, ...$). $\bar{\omega}_e$ is the fundamental vibrational frequency in $cm^{-1}$.When it is used: To find ideal, equally spaced vibrational energy levels.Common exam situations: Calculating the "Zero-Point Energy" by plugging in $\upsilon = 0$, proving that the molecule never stops vibrating even at absolute zero ($G(0) = \frac{1}{2}\bar{\omega}_e$).Anharmonic Vibrational Term ValueFormula: $G(\upsilon) = \bar{\omega}_e(\upsilon + \frac{1}{2}) - \bar{\omega}_e x_e(\upsilon + \frac{1}{2})^2 + \bar{\omega}_e y_e(\upsilon + \frac{1}{2})^3 ...$Meaning of each variable: $x_e$ and $y_e$ are anharmonicity constants. Note that usually $x_e \ll 1$.When it is used: When calculating the real-world spectrum where energy levels crowd closer together as the bond approaches the breaking point (dissociation).Common exam situations: You will be given the absorption frequencies for the fundamental band ($\upsilon=0 \to 1$) and the first overtone ($\upsilon=0 \to 2$) and asked to solve a system of equations to find the anharmonicity constant $x_e$ and the true frequency $\bar{\omega}_e$.Morse PotentialFormula: $E(R) = D_e [1 - e^{-a(R - R_e)}]^2$Meaning of each variable: $D_e$ is the well depth (dissociation energy). $a$ controls the width of the well.When it is used: It is the mathematical function used to graph the realistic, asymmetric "Anharmonic Oscillator" curve seen in Vibrational Spectra of Diatomic Molecule2.pdf.Common exam situations: Explaining why the harmonic oscillator fails at large displacements (because a parabola goes to infinity, meaning the bond never breaks), whereas the Morse potential tapers off to a finite dissociation energy $D_e$.

\end{document}